% fixed127 substantive insertion: reciprocity-marked Hilbert--Fox dictionary
% and the strict full-group marked Jacobian / pro-2 Fox chain comparison.
% This section is independent of all prior draft-file names and is included
% verbatim in the fixed127 master manuscript.

\subsection{Exact reciprocity marking and a strict dyadic marked-to-Fox bridge}
\label{subsec:dyadic-reciprocity-nielsen-bridge}

The abstract quadratic Fox form does not by itself fix a Kummer basis,
but the full reciprocity marking does.  On residual sectors with pro-$2$
image the marked tame direction also becomes contractible, leaving a
literal integral equivalence between the completed full-group
two-relation Jacobian and the genuine one-relator Fox resolution.
The resulting bridge is stronger than a comparison through the
intrinsic derived category, but is specific to the pro-$2$ residual
sectors of $G_{\mathbf Q_2}$.

\begin{theorem}[Exact reciprocity completion of the dyadic Hilbert--Fox basis]
\label{thm:dyadic-exact-hilbert-reciprocity-marking}
Fix the arithmetic local reciprocity map
$\operatorname{rec}:\mathbf Q_2^\times\to
G_{\mathbf Q_2}^{\mathrm{ab}}$,
so that $\operatorname{rec}(2)$ is arithmetic
Frobenius.  In the Roe--Turturean fully marked
Demu\v{s}kin identification
\cite{RoeTurtureanGQ2}, the pro-$2$
abelianizations of the standard generators satisfy
\begin{equation}\label{eq:dyadic-reciprocity-generator-images}
\boxed{\quad
\bar a=\operatorname{rec}(-4),\qquad
\bar s=\operatorname{rec}(1/2),\qquad
\bar y=\operatorname{rec}(-3).
\quad}
\end{equation}
Consequently the Hilbert--Fox relation
\eqref{eq:q2-kummer-fox-markings}
has \emph{zero} residual shear:
\begin{equation}\label{eq:dyadic-exact-kummer-fox}
\boxed{\quad
\kappa_{-1}=\chi_a,\qquad
\kappa_5=\chi_s,\qquad
\kappa_2=\chi_y.
\quad}
\end{equation}
More explicitly, for the ordered generators $(a,s,y)$
and ordered square classes $(-1,5,2)$, the matrix
of quadratic-character evaluations is exactly $I_3$:
\begin{equation}\label{eq:dyadic-exact-hilbert-evaluation}
\bigl(\kappa_u(g)\bigr)_{
u\in(-1,5,2),\ g\in(a,s,y)}
=
\begin{pmatrix}1&0&0\\0&1&0\\0&0&1\end{pmatrix},
\end{equation}
where $1\in\mathbf F_2$ means Hilbert symbol $-1$.
Writing a square class as
$u=(-1)^b5^c2^h$ with $b,c,h\in\mathbf F_2$,
the Hilbert pairing between
$u=(-1)^b5^c2^h$ and
$v=(-1)^{b'}5^{c'}2^{h'}$ becomes
\begin{equation}\label{eq:dyadic-square-class-hilbert-polar}
\boxed{\quad
(u,v)_2=(-1)^{\,bb'+ch'+hc'}.
\quad}
\end{equation}
This is exactly the bilinear pairing supplied by the
Fox degree-two symbol $A^2+SY+YS$.
\end{theorem}

\begin{proof}
The reciprocity identities
\eqref{eq:dyadic-reciprocity-generator-images}
are the marked abelianization of
Roe--Turturean, Lemma~3.5
\cite{RoeTurtureanGQ2}, and supply precisely
the information not read from the abstract cup matrix.
For a Kummer character, the evaluation on
$\operatorname{rec}(v)$ is the sign of the
quadratic norm-residue Hilbert symbol $(u,v)_2$;
using the opposite reciprocity normalization
would invert $\operatorname{rec}(v)$, but has
no effect on a quadratic character.
Apply Equation~\eqref{eq:q2-hilbert-explicit}
to the pairs
\[
 (u,v),\qquad
 u\in\{-1,5,2\},\quad
 v\in\{-4,1/2,-3\}.
\]
For $u=-1$, the symbols are $(-1,+1,+1)$;
for $u=5$, they are $(+1,-1,+1)$;
and for $u=2$, they are $(+1,+1,-1)$.
In particular,
$(2,1/2)_2=(2,2)_2=+1$, which fixes the
formerly undetermined value
$\kappa_2(s)=0$.
This gives \eqref{eq:dyadic-exact-hilbert-evaluation}
and \eqref{eq:dyadic-exact-kummer-fox}.
Finally the Hilbert pairing in the square-class
basis $(-1,5,2)$ has matrix
\eqref{eq:q2-hilbert-fox-matrix}.
Multiplying by the two coordinate column vectors
gives exactly
\eqref{eq:dyadic-square-class-hilbert-polar}.
\end{proof}

\begin{corollary}[Exact full-marked pro-$2$ Kummer coordinates]
\label{cor:dyadic-full-marked-kummer-coordinates}
Under the Nielsen transform
\begin{equation}\label{eq:dyadic-nielsen-full-generator-map}
\sigma=s,\qquad x_1=y,\qquad
x_0=s^{-2}a^{-1},
\qquad
a=x_0^{-1}\sigma^{-2},
\end{equation}
the pro-$2$ abelianized generator markings are
\begin{equation}\label{eq:dyadic-full-pro2-reciprocity}
\bar x_0=\operatorname{rec}(-1),\qquad
\bar\sigma=\operatorname{rec}(1/2),\qquad
\bar x_1=\operatorname{rec}(-3).
\end{equation}
Hence the three quadratic Kummer classes
are the dual basis of $(x_0,\sigma,x_1)$
in the order $(-1,5,2)$:
\begin{equation}\label{eq:dyadic-full-generator-evaluation}
\bigl(\kappa_u(g)\bigr)_{u=(-1,5,2),\,g=(x_0,\sigma,x_1)}
=I_3.
\end{equation}
In this ordering the mod-$2$ quadratic Fox
symbol of the full-marked pro-$2$ relation
is
\[
X_0^2+X_1\Sigma+\Sigma X_1.
\]
The $\nu_{\rm ur}$ and $\theta$ markings,
together with the fixed reciprocity classes,
therefore remove the last Hilbert--Fox
basis ambiguity without changing the
underlying cup form.
\end{corollary}

\begin{proof}
On abelianizations the identity
$x_0=s^{-2}a^{-1}$ gives
\[
 \bar x_0=
 \operatorname{rec}\bigl((1/2)^{-2}(-4)^{-1}\bigr)
 =\operatorname{rec}(-1).
\]
The two other identities follow immediately from
\eqref{eq:dyadic-reciprocity-generator-images}.
The Hilbert computations in
Theorem~\ref{thm:dyadic-exact-hilbert-reciprocity-marking}
then prove \eqref{eq:dyadic-full-generator-evaluation}.
For the quadratic relation, in the notation
$P_0=x_0^{\sigma^2}x_0$ and
$Q_0=[x_1,\sigma]$, its degree-two Magnus terms
are $X_0^2$ and $X_1\Sigma+\Sigma X_1$:
conjugation by $\sigma^2$ changes the
degree-one Magnus component of $x_0$
only in degree at least three.
\end{proof}

\begin{theorem}[Pro-$2$ residual compression of the marked full group]
\label{thm:dyadic-full-marked-pro2-compression}
Let $E/\mathbf Q_2$ be finite,
$\bar\rho:G_{\mathbf Q_2}\to\mathrm{GL}_r(k_E)$
be a genuine residual representation
whose finite image is a $2$-group,
and let $A$ be a complete Noetherian
local $\mathcal O_E$-algebra with residue
field $k_E$.  Every framed continuous
lift $\rho_A$ of $\bar\rho$ has pro-$2$
image and factors through
$P=G_{\mathbf Q_2}(2)$.
In the Roe--Turturean full marked presentation
\cite{RoeTurtureanGQ2}, the identities
\begin{equation}\label{eq:dyadic-pro2-marked-specialization}
\rho_A(\tau)=I_r,\qquad
\rho_A(\sigma^{\omega_2})=\rho_A(\sigma),
\qquad
\rho_A((x_i\tau)^{\omega_2})=\rho_A(x_i)
\quad (i=0,1)
\end{equation}
hold.  The tame relator is thereby eliminated,
and the full marked wild relator becomes
\begin{equation}\label{eq:dyadic-full-pro2-relator}
\boxed{\quad
R=P_0Q_0,\quad
P_0=x_0^{\sigma^2}x_0,\quad
Q_0=[x_1,\sigma].
\quad}
\end{equation}
This gives an equality of framed deformation
functors and therefore an isomorphism of
complete local representing rings
\begin{equation}\label{eq:dyadic-pro2-deformation-ring-equality}
R^\square_{\bar\rho}(G_{\mathbf Q_2})
\xrightarrow{\;\sim\;}
R^\square_{\bar\rho}(P).
\end{equation}
The two full-group marked equations on
this residual sector are ordinary
finite-word matrix equations once the
canonical pro-$2$ specialization is made.
No infinite profinite binomial series is
needed \emph{on this sector}; outside
pro-$2$ residual sectors, the formula of
Theorem~\ref{thm:dyadic-omega2-power-formula}
is still necessary.
\end{theorem}

\begin{proof}
Reduction identifies the image of $\rho_A$
modulo $\mathfrak m_A$ with the finite
$2$-group $\operatorname{im}\bar\rho$.
The congruence kernel of
$\mathrm{GL}_r(A)\to\mathrm{GL}_r(k_E)$
is pro-$2$ for the maximal-ideal
topology.  Therefore the compact image
$\rho_A(G_{\mathbf Q_2})$ is an extension
of pro-$2$ groups and is pro-$2$.
The universal property of the maximal
pro-$2$ quotient proves factorization.

In a pro-$2$ group, a relation
$g^{-1}tg=t^2$ forces $t=1$:
in any finite $2$-group quotient, a
nontrivial element cannot be conjugate
to its square since the two have
different orders.  Apply this to the
tame relation
$\tau^\sigma=\tau^2$.
The profinite idempotent $\omega_2$
acts as the identity on all pro-$2$
elements, so
$\sigma_2=\sigma$ and
$u_i=(x_i\tau)^{\omega_2}=x_i$.
The remaining marked words then simplify
successively as
$d_0=c_0=d_g=h_c=1$,
$g_0=\sigma^2$, and
$h_0=x_0^{\sigma^2}x_0$.
The full wild relation
$h_0u_1^{-1}x_1^\sigma c_0=1$
is precisely
\eqref{eq:dyadic-full-pro2-relator}.
The marked condition on the normal closure
of $x_0,x_1$ is automatically satisfied
in every pro-$2$ quotient.
Conversely every representation of the
one-relator pro-$2$ group in this residual
chart defines an admissible full-group
marked representation by putting
$\tau=1$.  Hence the framed lift functors
are naturally isomorphic on every
Artinian coefficient algebra; continuity
and completeness give
\eqref{eq:dyadic-pro2-deformation-ring-equality}.
The last assertion follows since all
remaining operations in $P_0Q_0$
are ordinary finite words.
\end{proof}

\begin{theorem}[Strict integral Nielsen--Fox chain equivalence]
\label{thm:dyadic-strict-nielsen-fox-chain}
Retain Theorem~\ref{thm:dyadic-full-marked-pro2-compression},
let $B$ be a complete local coefficient algebra,
and let $M$ be a finite projective $B$-module
with continuous $P$-action.
Set
\[
S=\rho_M(s),\quad A_0=\rho_M(a),\quad
X_0=\rho_M(x_0)=S^{-2}A_0^{-1},
\quad
\Sigma=\rho_M(\sigma)=S,
\]
and let $Q_0=[x_1,\sigma]$ as above.
For a degree-one standard Fox coordinate
$(v_a,v_s,v_y)\in M^3$ define
\begin{equation}\label{eq:dyadic-strict-nielsen-s1}
\mathsf N^1(v_a,v_s,v_y)=
\left(
 -S^{-2}A_0^{-1}v_a
 -(S^{-1}+S^{-2})v_s,\quad
 v_s,\quad v_y
\right),
\end{equation}
where the target is in $(x_0,\sigma,x_1)$
coordinate order.  Put
\begin{equation}\label{eq:dyadic-strict-nielsen-s2}
\mathsf N^0=\operatorname{id}_M,\qquad
\mathsf N^2=-\rho_M(Q_0)^{-1}.
\end{equation}
Then $\mathsf N^\bullet$ is an \emph{invertible
strict integral chain map} from the
standard one-relator Fox complex of
$a^2s^4[s,y]$ to the marked pro-$2$
Fox complex of $P_0Q_0$:
\begin{equation}\label{eq:dyadic-strict-nielsen-chain-square}
\mathsf N^1 d_{\rm std}^0
=d_{\rm mark}^0\mathsf N^0,
\qquad
d_{\rm mark}^1\mathsf N^1
=\mathsf N^2d_{\rm std}^1.
\end{equation}
Its degree-one inverse is explicit:
\begin{equation}\label{eq:dyadic-strict-nielsen-inverse}
\begin{aligned}
v_s&=v_\sigma,\qquad v_y=v_{x_1},\\
v_a&=-X_0^{-1}v_{x_0}
     -X_0^{-1}(\Sigma^{-1}+\Sigma^{-2})v_\sigma.
\end{aligned}
\end{equation}
The chain isomorphism commutes strictly with
integral coefficient base change and
preserves the intrinsic cup pairing
on cohomology.  The chosen Fox diagonals
need not be strictly identified as
multiplicative chain maps.
\end{theorem}

\begin{proof}
For the left crossed-derivation convention,
$c(gh)=c(g)+\rho_M(g)c(h)$ and
$c(g^{-1})=-\rho_M(g)^{-1}c(g)$.
By the Nielsen identity
$x_0=s^{-2}a^{-1}$, direct evaluation gives
\[
c(x_0)=
 -(S^{-1}+S^{-2})c(s)
 -S^{-2}A_0^{-1}c(a),
\]
which is exactly
\eqref{eq:dyadic-strict-nielsen-s1}.
The inverse word
$a=x_0^{-1}\sigma^{-2}$
gives \eqref{eq:dyadic-strict-nielsen-inverse}.
For a principal degree-zero coboundary
$c(g)=(\rho_M(g)-1)v$,
both sides evaluate on the same group
elements, proving the first identity of
\eqref{eq:dyadic-strict-nielsen-chain-square}.

The relation words satisfy the exact
\emph{free-group} identity
\begin{equation}\label{eq:dyadic-nielsen-relator-conjugacy}
r_{\rm std}=a^2s^4[s,y]
          =(Q_0P_0)^{-1},\qquad
r_{\rm mark}=P_0Q_0
          =Q_0^{-1}r_{\rm std}^{-1}Q_0
\end{equation}
after replacing $a=x_0^{-1}\sigma^{-2}$,
$s=\sigma$, and $y=x_1$.
Because $\rho_M(r_{\rm std})=1$,
the crossed-derivation identity yields
\[
c(r_{\rm mark})
=-\rho_M(Q_0)^{-1}c(r_{\rm std}).
\]
Evaluation of arbitrary generator
$1$-cochains on the corresponding
free relator is precisely the completed
Fox differential $d^1$.
This proves the second identity of
\eqref{eq:dyadic-strict-nielsen-chain-square}
without requiring the generator values
to be cocycles.
Both $\mathsf N^1$ and $\mathsf N^2$
are invertible integral matrices
over $B$ and involve only transport
matrices and their inverses; therefore
base change is strict.
The invariant cup product is that of
$R\Gamma(P,M)$ (or its paired coefficient
objects); identifying both Fox complexes
with the same intrinsic cochains
shows that cup classes agree.
\end{proof}

\begin{theorem}[Tame-block contraction of the full marked Jacobian]
\label{thm:dyadic-tame-jacobian-fox-splitting}
Let $B$ and $M$ be as in
Theorem~\ref{thm:dyadic-strict-nielsen-fox-chain},
with the coefficient action factoring through $P$.
On the genuine pro-$2$ residual chart,
let $\mathcal J^\bullet_{\rm full}(M)$ be
the three-term \emph{linearized marked
relation complex}, with
\[
\mathcal J^0=M,\qquad
\mathcal J^1=M^4
\ \text{in coordinates }(\sigma,\tau,x_0,x_1),
\qquad
\mathcal J^2=M^2
\ \text{in the tame/wild relation order}.
\]
The degree-one differential evaluates
left crossed derivations on the
two full marked relators, and the
degree-zero differential is the
usual principal-cocycle map.
Writing $\Sigma=\rho_M(\sigma)$,
put
\begin{equation}\label{eq:dyadic-tame-L-unit}
\mathsf L=\Sigma^{-1}-2I_M
\in\operatorname{Aut}_B(M).
\end{equation}
After reordering degree-one coordinates as
$(x_0,\sigma,x_1;\tau)$, the marked
relation differential has block matrix
\begin{equation}\label{eq:dyadic-full-jacobian-block}
d^1_{\rm full}
=
\begin{pmatrix}
 0&\mathsf L\\
 d^1_{\rm mark}&\mathsf B_\tau
\end{pmatrix},
\end{equation}
where $\mathsf B_\tau:M\to M$ is the
finite-word derivative of the wild
relation in the tame direction at
$\tau=1$.  It also has the following
\emph{closed transport formula}.  Put
\[
X=\rho_M(x_0),\quad
G=\Sigma^2,\quad
Z=\Sigma^{-1}X\Sigma,\quad
P_0^\rho=G^{-1}XGX.
\]
Then
\begin{equation}\label{eq:dyadic-wild-tau-jacobian-closed}
\boxed{\quad
\mathsf B_\tau
=P_0^\rho(G^{-1}X+3X-I_M)
 +(Z^{-1}-I_M)X.
\quad}
\end{equation}
The output change
\begin{equation}\label{eq:dyadic-tame-row-elimination}
(u_{\rm tame},u_{\rm wild})
\longmapsto
\bigl(u_{\rm tame},
u_{\rm wild}-\mathsf B_\tau\mathsf L^{-1}u_{\rm tame}\bigr)
\end{equation}
is integral and invertible.
It yields an actual chain isomorphism
\begin{equation}\label{eq:dyadic-full-marked-chain-splitting}
\boxed{\quad
\mathcal J^\bullet_{\rm full}(M)
\cong
C^\bullet_{\rm Fox}(P,M)
\oplus
\bigl(0\to M
  \xrightarrow{\ \mathsf L\ }M\to0\bigr),
\quad}
\end{equation}
with the acyclic summand supported in
cohomological degrees $1$ and $2$.
Composing with
Theorem~\ref{thm:dyadic-strict-nielsen-fox-chain}
gives a \emph{direct integral chain-level
comparison} between the marked
full-$G_{\mathbf Q_2}$ Jacobian and the
standard dyadic Demu\v{s}kin Fox--Spencer
coefficient complex.

In this pro-$2$ image sector the full
two-relator marked Jacobian is therefore
cohomologically Fox-complete in degrees
$0,1,2$ and computes
$R\Gamma(G_{\mathbf Q_2},M)$.
This conclusion is \emph{not} extended to
arbitrary residual images, for which
the fixed full-group marked relator
Jacobian may retain higher syzygies
and the all-Sylow derived construction
remains necessary.
\end{theorem}

\begin{proof}
At a genuine pro-$2$ representation the
tame generator $\tau$ acts as the identity,
and by Theorem~\ref{thm:dyadic-full-marked-pro2-compression}
every occurrence of $\omega_2$ in the
marked relations is the identity
operation on the ambient pro-$2$
deformation sector, including first-order
semidirect extensions by $M$.
The tame relation can be written as
\[
 r_t=\sigma^{-1}\tau\sigma\tau^{-2}=1.
\]
For an arbitrary crossed derivation $c$
with $c(\sigma)=v_\sigma$ and
$c(\tau)=v_\tau$, and
$\rho_M(\tau)=I$, one calculates
\[
 c(r_t)=(\Sigma^{-1}-2I_M)v_\tau;
\]
all $v_\sigma$ terms cancel.
Modulo the maximal ideal of $B$,
the reduction of $\mathsf L$ is
$\bar\Sigma^{-1}$ since $2=0$ there,
hence $\mathsf L$ is invertible over
the complete local ring $B$.
For completeness, the closed
$\tau$-derivative in
\eqref{eq:dyadic-wild-tau-jacobian-closed}
is obtained directly from the literal
Roe--Turturean auxiliary words.
At $\tau=1$ and with $c(\tau)=v_\tau$
and all other generator variations zero,
\[
c(d_0)=Xv_\tau,\qquad
c(d_g)=G^{-1}Xv_\tau,\qquad
c(h_c)=0,\qquad
c(c_0)=(Z^{-1}-I_M)Xv_\tau.
\]
Since $h_0=(x_0^{g_0}x_0)d_gd_0^3h_c$
and $u_1=x_1\tau$, the first two
contributions to $c(r_{\rm wild})$ are
$P_0^\rho(G^{-1}X+3X)v_\tau$
and $-P_0^\rho v_\tau$, followed by
the displayed contribution from $c_0$.
This proves the closed formula
without an unspecified Fox derivative.

The derivative of the wild relation
with $\tau$ fixed equal to $1$ is
the Fox derivative of $P_0Q_0$;
its derivative with respect to
$v_\tau$ is by definition
$\mathsf B_\tau$.  This proves the
block display
\eqref{eq:dyadic-full-jacobian-block}.
The row operation
\eqref{eq:dyadic-tame-row-elimination}
kills $\mathsf B_\tau$ without changing
the first Fox block.

Since $\tau$ acts trivially on $M$,
the degree-zero principal cocycle has
$v_\tau=(\rho_M(\tau)-1)v=0$;
hence the row/column decomposition
also splits the entire complex,
not only its degree-one differential.
The one-block complex with differential
$\mathsf L$ is contractible, yielding
\eqref{eq:dyadic-full-marked-chain-splitting}.
The Nielsen chain maps above then give
the explicit comparison in the standard
$a,s,y$ coordinates.
The standard Fox resolution of the
Demu\v{s}kin group $P$ computes
$R\Gamma(P,M)$ by
Theorem~\ref{thm:fox-spencer-demushkin},
and inflation from $P$ computes
$R\Gamma(G_{\mathbf Q_2},M)$ for
pro-$2$ coefficient actions by
Theorem~\ref{thm:sylow-fox-retract}.
No claim about the naive marked complex
on other residual sectors is used.
\end{proof}

\begin{corollary}[Exact full-group rank-one hypersurface elimination]
\label{cor:dyadic-full-rankone-tame-elimination}
For the trivial residual rank-one
representation $\bar\rho=\mathbf1$,
write the four marked scalar generator
matrices as
\[
\Sigma=1+z_\sigma,\quad
T=1+z_\tau,\quad
X_0=1+z_0,\quad
X_1=1+z_1.
\]
The completed full-group tame and wild
marked relation ideal is
\begin{equation}\label{eq:dyadic-full-rankone-two-relator-ideal}
\boxed{\quad
I_{\rm marked}
=(z_\tau,\ z_0(z_0+2))
\subseteq
\mathcal O_E[[z_\sigma,z_\tau,z_0,z_1]].
\quad}
\end{equation}
Thus the full-group marked
complete local presentation is already
a regular sequence, with a contractible
tame variable and one genuine
dyadic hypersurface:
\[
R^\square_{\mathbf1}
\cong
\mathcal O_E[[z_\sigma,z_0,z_1]]/
(z_0(z_0+2)).
\]
Under the Nielsen coordinate transform
\begin{equation}\label{eq:dyadic-rankone-u-z0-change}
 1+u=X_0^{-1},\qquad
 u=-\frac{z_0}{1+z_0},
\end{equation}
the identity of ideals is
\begin{equation}\label{eq:dyadic-rankone-hypersurface-comparison}
u(u+2)=-\frac{z_0(z_0+2)}{(1+z_0)^2}.
\end{equation}
Consequently the two full-group marked
relations and the one-relator
Demu\v{s}kin deformation hypersurface
of Theorem~\ref{thm:dyadic-q2-rankone-integral-ring}
agree \emph{integrally and scheme-theoretically},
not merely on their tangent spaces
or after inverting $2$.
\end{corollary}

\begin{proof}
Every scalar lift of $\mathbf1$ takes
values in the pro-$2$ principal units.
Conjugation is trivial, and the tame
relation $\tau^\sigma=\tau^2$
therefore becomes $T=T^2$,
equivalently $z_\tau=0$ since
$T\in1+\mathfrak m_A$ is invertible.
All $\omega_2$ powers act as the identity
on these principal units.  More strongly,
for arbitrary scalar $T$ near $1$ the
\emph{literal} full marked tame and wild
relations simplify exactly to
\[
r_t=T^{-1},\qquad
r_{\rm wild}=X_0^2T^3.
\]
Indeed $d_0=d_g=T$, $h_c=c_0=1$,
$h_0=X_0^2T^4$, and
$u_1^{-1}x_1^\sigma=T^{-1}$.
With $T=1$, the wild relation
reduces to $P_0Q_0$ and, for scalar
commuting matrices, to $X_0^2=1$.
Thus its equation is exactly
$z_0(z_0+2)=0$ modulo $z_\tau$.
Elementary ideal elimination gives
\eqref{eq:dyadic-full-rankone-two-relator-ideal}
in the complete local ambient
power-series ring; the tame relation
has the explicit unit multiple
$T^{-2}(T-T^2)$ of $z_\tau$,
and the wild relation differs from
$X_0^2-1$ by an element of
$(z_\tau)$.  Both generators form
a regular sequence.
Finally
$a=x_0^{-1}\sigma^{-2}$ and $s=\sigma$
give $a s^2=x_0^{-1}$ in the
scalar quotient, so
$1+u=\rho(a)\rho(s)^2=X_0^{-1}$.
The displayed rational power-series
identity
\eqref{eq:dyadic-rankone-hypersurface-comparison}
is integral because $1+z_0$ is a
unit in the completed local ring.
This proves the exact identification.
\end{proof}

\begin{remark}[Sharp scope of the direct integral bridge]
\label{rem:dyadic-full-chain-bridge-scope}
If $\operatorname{im}\bar\rho$ is a
$2$-group, its Sylow preimage
in Theorem~\ref{thm:sylow-fox-retract}
is all of $G_{\mathbf Q_2}$.
The present chain equivalence therefore
makes the pointwise Sylow--Fox
comparison completely native on
these sectors, without choosing
a chain homotopy inverse of a raw
Herr representative.
For a general residual image,
the Sylow preimage is a proper
finite-index subgroup; the
full-group marked relation
Jacobian need not be cohomologically
complete, and an explicit
degree-two integral comparison
requires actual Schreier
relations and syzygy control.
The all-Sylow/derived construction
retains the unconditional result.
Neither this restriction nor the
finite-jet profinite-power formula
of fixed126 is a claim of a uniform
four-generator full-group
cohomological Fox resolution
for all residual sectors.
\end{remark}
